\section{Introducción}

La respuesta macroscópica de un material frente a un campo magnético externo se describe mediante la magnetización $M$, definida como el momento magnético por unidad de volumen. Para materiales en los que las interacciones entre dipolos magnéticos son débiles, la magnetización resulta proporcional al campo magnetizante aplicado $H$,

\begin{equation}
M=\chi H,
\label{MchiH}
\end{equation}

donde $\chi$ es la susceptibilidad magnética del material. Esta relación lineal permite clasificar a los materiales según el signo y magnitud de $\chi$. \cite{Bruvera}

Los materiales diamagnéticos presentan susceptibilidad negativa ($\chi<0$). En ellos, el campo aplicado induce corrientes que generan un momento magnético opuesto al campo externo, produciendo una débil repulsión magnética. Este comportamiento es característico de materiales cuyos átomos poseen capas electrónicas completas y, por lo tanto, momento magnético neto nulo, como los gases inertes. Por el contrario, los materiales paramagnéticos poseen momentos magnéticos permanentes asociados a electrones desapareados que tienden a alinearse con el campo aplicado, produciendo una susceptibilidad positiva ($\chi>0$). Debido a la agitación térmica, dicha alineación es sólo parcial y la susceptibilidad disminuye al aumentar la temperatura. \cite{Bruvera} \cite{Kittel} 

El origen microscópico de estos fenómenos se encuentra en los momentos magnéticos asociados al movimiento orbital y al $spin$ de los electrones. Cada electrón puede interpretarse como un pequeño dipolo magnético caracterizado por un momento magnético $\mu$ capaz de interactuar tanto con campos externos como con otros momentos magnéticos presentes en el material. En el caso más simple de partículas con $spin$ $(1/2)$, el sistema puede modelarse mediante dos niveles energéticos, correspondientes a las orientaciones paralela y antiparalela respecto de un campo magnético aplicado y cuya ocupación relativa depende de la temperatura según la distribución de Boltzmann \cite{Kittel} 

También existen materiales que, por su composición, no cumplen con \eqref{MchiH}. En estos materiales $M(H)$ no sólo no es lineal, sino que ni siquiera es monovaluada y su valor depende de la historia del sistema. Estas interacciones, conocidas como interacciones de intercambio, tienen origen cuántico y son responsables de la aparición de orden magnético colectivo. Dependiendo de la orientación relativa favorecida entre momentos adyacentes, los materiales pueden clasificarse como antiferromagnéticos, ferrimagnéticos o ferromagnéticos. En los antiferromagnetos los momentos se alinean antiparalelamente dando lugar a una magnetización neta nula. Los ferrimagnetos presentan también subredes antiparalelas, aunque con momentos de distinta magnitud, por lo que conservan una magnetización macroscópica distinta de cero. Finalmente, en los ferromagnetos la interacción de intercambio favorece la alineación paralela de los momentos magnéticos, originando una magnetización espontánea incluso en ausencia de campo externo, como es el caso de la mayoria de los materiales ordinarios. \cite{Bruvera}

Mientras que las interacciones de intercambio favorecen a la alineación de los espines, la agitación térmica tiende a desordenarlos. A temperaturas suficientemente elevadas, el efecto de la temperatura domina sobre las interacciones de intercambio y destruye el el orden magnético de largo alcance. Como consecuencia, existe una temperatura de Curie $T_C$, que separa la fase ferromagnética ordenada de la fase paramagnética desordenada. Por debajo de $T_C$ el sistema presenta magnetización espontánea, mientras que por encima de ella dicha magnetización desaparece. Desde el punto de vista de la física estadística, esta transición constituye un ejemplo de ruptura espontánea de simetría, ya que el sistema selecciona una dirección preferencial de magnetización aun cuando no se distingue ninguna dirección particular del espacio. \cite{Kittel} \cite{Aharoni}

En muestras macroscópicas, el estado ferromagnético suele organizarse en dominios magnéticos, regiones donde los momentos permanecen aproximadamente alineados. Esta configuración reduce la energía magnetostática del sistema y explica por qué muchos materiales ferromagnéticos presentan una magnetización neta nula en ausencia de campo externo. Cuando se aplica un campo magnético, la configuración de dominios se modifica para favorecer aquellas orientaciones más próximas a la dirección del campo. Sin embargo, este proceso no es completamente reversible: la respuesta del material depende del estado magnético previo y de la trayectoria seguida por el campo aplicado. Como consecuencia, la magnetización no queda determinada únicamente por el valor instantáneo del campo, dando lugar al fenómeno de histéresis. Los ciclos de histéresis se caracterizan por una magnetización remanente $M_r$, una magnetización de saturación $M_s$ y un campo coercitivo $H_c$, definido como el campo necesario para anular la magnetización remanente. \cite{Bruvera}
 
La dependencia de la magnetización con la temperatura presenta distintos comportamientos según el rango considerado. A temperaturas suficientemente bajas respecto de $T_C$, la disminución de la magnetización se debe principalmente a excitaciones colectivas del $spin$ denominadas magnones y puede describirse mediante la ley de Bloch

\begin{equation}
M(T)=1-\left(\frac{T}{T_C}\right)^{3/2}
\label{eq:Bloch}
\end{equation}

Por otro lado, al aproximarse a la temperatura de Curie, las fluctuaciones críticas dominan el comportamiento del sistema y la magnetización espontánea sigue una ley de potencia derivada del modelo propuesto por Heisenberg

\begin{equation}
M(T)=M(0)\left(1-\frac{T}{T_C}\right)^\beta,
\label{Heisemberg}
\end{equation}

donde $M(0)$ es la magnetización supuesta para $T=0$ $K$ y $\beta$ es el exponente crítico asociado a la transición ferromagnética. En la aproximación desarrollada por Weiss y Curie se obtiene $\beta=1/2$, mientras que los valores experimentales observados en numerosos materiales suelen encontrarse en el intervalo $0.33-0.37$. \cite{Aharoni} \cite{Lue}

En este trabajo se estudiará experimentalmente la dependencia de la magnetización con la temperatura en una ferrita ..., analizando tanto el régimen descrito por la ley de Bloch como el comportamiento de la magnetización en las proximidades de la temperatura de Curie. A partir de estas mediciones se determinarán los parámetros característicos del material y buscara el ratio de Mn/Zn de la muestra.
